%
%
% wabun2025.tex
%
%
%
\documentclass[fontsize=11pt]{jlreq}
\usepackage{subfiles}
\usepackage{amsmath, amssymb, braket} % AMSの数式
\usepackage{lmodern} % 美文書入門 p.74 5.2数式用のフォント
\usepackage{graphicx,color} % https://texwiki.texjp.org/?hyperref
\usepackage{subcaption} % figureなどのキャプション
\usepackage{bm} % bold math
\usepackage{quantikz} % 量子回路図
\usepackage[version=4]{mhchem} % 化学式
\usepackage[backend=biber, style=numeric, maxnames=5, doi=false, url=true, eprint=true]{biblatex}
\addbibresource{quantum.bib}
\makeatletter
\let\mkbibemph=\mkbibitalic
\let\blx@imc@mkbibemph=\blx@imc@mkbibitalic
\makeatother
\setlength{\biblabelsep}{0.5em}

%\usepackage{pdfx}  % side-effect: hyperrefの設定もされる
\usepackage{hyperref} % リンクの埋め込み、pdfファイルの目次の生成
\hypersetup{
    colorlinks=true,
    linkcolor=blue,
    urlcolor=cyan,
    pdftitle={Quantum Circuit for First Quantization Hamiltonian Simulation Handling Coulomb Potential Pole},
    pdfauthor={Hideo Takahashi},
    pdfsubject={Research Paper},
    pdfkeywords={Quantum Computer, First Quantization, ab initio, Simulator, Coulomb Potential, Quantum Chemistry},
    bookmarksnumbered=true,
    bookmarksopen=true
}
\usepackage{xr}
\externaldocument{wabun2025_sm}
\usepackage{paper2025}

\begin{document}
\title
{フル量子回路による1,2次元水素原子モデルの\\
時間発展シミュレーション:誤差解析とパラメータ最適化}
\author{高橋 英男}
\date{\today}
\maketitle
\begin{abstract}
本研究では、30量子ビット未満の小規模な量子コンピュータエミュレータ上で動作可能
な、第一量子化に基づくハミルトニアンシミュレータの量子回路を構築した。
具体的には、26量子ビット（座標レジスタ6ビット）を用い、単一水素原子の2次元モデ
ルのシミュレーションを実現した。我々の知る限り、これはクーロンポテンシャルを含め
全て量子回路で実装したグリッド法シミュレータの初の動作報告である。

限られた量子リソース(qubit数)で回路を構築するため、クーロンポテンシャルの計算に
はUniformly Controlled Rotation (UCR) 回路を適用した。さらに、高精度な計算結果を
得るための理論的指針として、クーロンポテンシャルの特異点（極）の扱いと
、離散化パラメータの最適化手法を確立した。
クーロンポテンシャルの特異点については、極を中心とする微小な円形領域における
ポテンシャルの単純な幾何学的平均から実効値を定めることで、発散を回避しつつ誤差を
最小化できることを示した。

また、与えられたqubit数の制約下において、空間グリッド間隔は「空間離散化誤差」と
「領域打切り誤差」のトレードオフ評価から最適値を決定する手法を提案した。時間発展
の刻み幅については、最悪ケースを想定した形式的なTrotter誤差の上限ではなく、波動
関数の最大周波数成分に着目したサンプリング定理に基づくことで、量子回路の無駄な深
さ（Depth）の増大を避けつつ十分な精度を保つ条件を明らかにした。本手法は、限られ
たリソース下で実用的な量子シミュレーションを設計するための強力な指針となる。
\end{abstract}

\subfile{wabun2025_1_intro.tex}
\subfile{wabun2025_2_method.tex}
\subfile{wabun2025_3_results.tex}
\subfile{wabun2025_4_analysis.tex}

\section{結論}
本研究では、2次元水素原子モデルを対象に、Qiskitエミュレータ上で全系の時間発展が
可能な回路構成を示した。我々の知る限り、これはクーロンポテンシャルを含むグリッド
法に基づく時間発展計算を量子回路として実装し、エミュレータ上で動作を確認した初の
報告である。

クーロンポテンシャルの特異点の扱いに関して、特異点を含む格子点における有効
ポテンシャルを定める実効半径として、2次元モデルにおいては$\reff=\deltar/4$を
用いる手法を提案した。数値実験によりこの値が系のエネルギー誤差を最小化することが
確認されただけでなく、これが特異点周辺の幾何学的な領域平均から導出される値と
一致することを示し、離散化空間における本手法の物理的な妥当性を明らかにした。

さらに、与えられたqubit数（$\nb$）の下で誤差を最小化するシミュレーションパラメータ
の決定手順を確立した。空間グリッド間隔$\deltar$は、空間離散化誤差と領域打切り誤差
のトレードオフから最適値およびシミュレーションのダイナミックレンジとして決定できる
ことを示した。また、Trotter分解の時間刻み$\deltat$については、最悪ケースを想定
した形式的なTrotter誤差の上限ではなく、波動関数の最大周波数成分に基づくサンプリング
定理の要請（半周期条件）から決定できることを数値的に確認した。特に対象が基底状態
付近の場合、この条件を満たす範囲において時間誤差よりも空間誤差が支配的になる（誤差が
飽和する）ことを示し、無駄な量子回路の深さ（Depth）の増大を回避するための指針を与えた。

これらの知見は、限られた量子リソースのもとでグリッド法を用いた量子シミュレーション
を設計・実行する上で、極めて実用的かつ強力な指針となることが期待される。

\printbibliography[heading=bibintoc, title={参考文献}]

%%%%%%%%%%%%%%%% 付録

\appendix
\section{局所平均ポテンシャルの実効半径の導出}
\label{chap:appendix_derivation_derivation}

第\ref{sec:coulomb-potential-handling}節で述べた近似式 $\Hla(r)$ におけるパラメータ $\reff$
の導出を示す。

まず、$\PsiTwoD_{0,0}$に関するパラメタ$\Delta_{1,0,0}$についての式\eqref{eq:delta-1-0-0}の導出を示す。

積分 $W_{0,0}$ は、
\begin{align}
  W_{0,0}
   &= \int_{r=0}^{\deltar/2}\int_{\theta=0}^{2\pi}
      \Hep(r)\abs{\PsiTwoD_{0,0}(r)}^2rd\theta dr\notag\\
   &= \int_{r=0}^{\deltar/2}\int_{\theta=0}^{2\pi}
      -\frac{1}{r}\left(\frac{q_0^3(0-\abs{0})!}{\pi(0+\abs{0})!}\right)
      (2q_0r)^{2\abs{0}}\exp(-2q_0r)
      \left[L_{0-\abs{0}}^{2\abs{0}}(2q_0r)\right]^2 rd\theta dr\notag\\
   &= -2\pi\int_{r=0}^{\deltar/2}\frac{q_0^3}{\pi}\frac{1}{r}e^{-2q_0r}r\ dr=-2q_0^3
      \left[-\frac{1}{2q_0}e^{-2q_0r}\right]_{0}^{\deltar/2}\notag\\
   &=-q_0^2(-e^{-q_0\deltar}+1)
\end{align}
となる。これに対して $W_{0,0}^\prime$ は、
\begin{align}
    W_{0,0}^\prime
    &=\pi\left(\frac{\deltar}{2}\right)^2
       \Hla(0)\abs{\PsiTwoD_{0,0}(0)}^2\notag\\
    &=\frac{\pi\deltar^2}{4}\left(-\frac{1}{\reff}\right)
    \left(\frac{q_0^3(0-\abs{0})!}{\pi(0+\abs{0})!}\right)
    (0)^{2\abs{0}}\exp(0)\left[L_{0-\abs{0}}^{2\abs{0}}(0)\right]^2\notag\\
    &= -\frac{\pi\deltar^2q_0^3}{4\pi\reff}\notag\\
    &= -\frac{\deltar^2q_0^3}{4\reff}
\end{align}
となる。$W_{0,0}=W_{0,0}^\prime$を解くと、
\begin{equation}
-q_0^2\left(-e^{-q_0\deltar}+1\right)=-\frac{\deltar^2q_0^3}{4\reff}
\end{equation}
\begin{align}
\reff&=\frac{\deltar^2q_0}{4\left(-e^{-q_0\deltar}+1\right)}\notag\\
&\simeq\frac{\deltar^2q_0}{4\left(-(1-q_0\deltar)+1\right)}\notag\\
&=\frac{\deltar}{4}
\end{align}
となる。

次に、$\PsiTwoD_{1,0}$に関するパラメタ$\Delta_{1,1,0}$についての式\eqref{eq:delta-1-1-0}の導出を示す。

$W_{1,0}$と$W_{1,0}^\prime$をそれぞれ求めると、$W_{1,0}$は、
\begin{align}
  W_{1,0}
   &= \int_{r=0}^{\deltar/2}\int_{\theta=0}^{2\pi}
      \Hep(r)\abs{\PsiTwoD_{1,0}(r)}^2rd\theta dr\notag\\
   &= \int_{r=0}^{\deltar/2}\int_{\theta=0}^{2\pi}
      -\frac{1}{r}\left(\frac{q_0^3(1-\abs{0})!}{\pi(1+\abs{0})!}\right)
      (2q_0r)^{2\abs{0}}\exp(-2q_0r)
      \left[L_{1-\abs{0}}^{2\abs{0}}(2q_0r)\right]^2 rd\theta dr\notag\\
   &= -2\pi\int_{r=0}^{\deltar/2}\frac{q_0^3}{\pi}e^{-2q_0r}(1-2q_0r)^2dr\notag\\
   &= -2q_0^3\left[\left(-\frac{1}{2q_0}+2q_0r^2\right)e^{-2q_0r}\right]_{0}^{\deltar/2} \notag\\
   &= -2q_0^3\left\{\left(-\frac{1}{2q_0}+\frac{q_0\deltar^2}{2}\right)e^{-q_0\deltar}+\frac{1}{2q_0}\right\}\notag\\
   &= q_0^2\left\{\left(1-q_0^2\deltar^2\right)e^{-q_0\deltar}-1\right\}
\end{align}
となる。これに対して $W_{1,0}^\prime$ は、
\begin{align}
    W_{1,0}^\prime
    &=\pi\left(\frac{\deltar}{2}\right)^2
         \Hla(0)\abs{\PsiTwoD_{1,0}(0)}^2\notag\\
    &=\frac{\pi\deltar^2}{4}\left(-\frac{1}{\reff}\right)
    \left(\frac{q_0^3(1-\abs{0})!}{\pi(1+\abs{0})!}\right)
    (0)^{2\abs{0}}\exp(0)\left[L_{1-\abs{0}}^{2\abs{0}}(0)\right]^2\notag\\
    &= -\frac{\pi\deltar^2q_0^3}{4\pi\reff}\notag\\
    &= -\frac{\deltar^2q_0^3}{4\reff}
\end{align}
となる。$W_{1,0}=W_{1,0}^\prime$を解くと、
\begin{equation}
q_0^2\left\{\left(1-q_0^2\deltar^2\right)e^{-q_0\deltar}-1\right\}
=-\frac{\deltar^2q_0^3}{4\reff}
\end{equation}
\begin{align}
\reff&=\frac{\deltar^2q_0}{4\left\{\left(1-q_0^2\deltar^2\right)e^{-q_0\deltar}-1\right\}}\notag\\
&\simeq-\frac{\deltar^2q_0}{4\left\{\left(1-q_0^2\deltar^2\right)(1-q_0\deltar)-1\right\}}\notag\\
&=-\frac{\deltar^2q_0}{4\left(1-q_0^2\deltar^2-q_0\deltar+q_0^3\deltar^3-1\right)}\notag\\
&\simeq\frac{\deltar}{4}
\end{align}
が得られる。


\end{document}